\documentclass{article}

\usepackage[utf8]{inputenc}     % pdfLaTeX
\usepackage[T1]{fontenc}
\usepackage[magyar]{babel}
\usepackage{amsmath}

% Set page size and margins
% Replace `letterpaper' with `a4paper' for UK/EU standard size
\usepackage[a4paper,top=2cm,bottom=2cm,left=3cm,right=3cm,marginparwidth=1.75cm]{geometry}
%for arrows
\usepackage{textcomp}

\title{Igenevek}
\author{Lengyel Zoltán Péter\thanks{Csukás Ibolya tanítása alapján}}
\date{Legutóbb frissítve: 2026. Március 15.}

\begin{document}

\maketitle
\tableofcontents
\newpage

\begin{center}
    {\Huge Igenevek}
\end{center}
\begin{itemize}
    \item Igékből képzett szavak, amelyek egy másik szófaj jellegzetességét is hordozzák
\end{itemize}

\section{Főnévi igenév}
\begin{itemize}
    \item ige + -ni / -nni képző
    \item a mondatban főnévhez hasonló szerepet lát el \\
    $\begin{array}{l}
    \quad\text{\underline{Futni} jó dolog } \\
    \quad\quad\downarrow \\
    \text{A \underline{futás} jó dolog}
    \end{array}$
    \item A főnévi igenév ragozható (az egyértelműség kedvéért) \\
    \begin{tabular}{|l|l|} \hline
        E/1 & adnom \\ \hline
        E/2 & adnod \\ \hline
        E/3 & adnia \\  \hline
        T/1 & adnunk \\ \hline
        T/2 & adnotok \\ \hline
        T/3 & adniuk \\ \hline
    \end{tabular}
\end{itemize}

\section{Melléknévi igenév}
\subsection{Folyamatos melléknévi igenév}
\begin{itemize}
    \item ige + -ó / -ő képző
    \begin{itemize}
        \item pl. mosolygó, nevelő
    \end{itemize}
    \item Sok folyamatos melléknévi igenév szerepelhet főnévként is 
    \begin{itemize}
        \item pl. író, terítő, tanuló, költő, ebédlő
    \end{itemize}
\end{itemize}
\subsection{Befejezett melléknévi igenév}
\begin{itemize}
    \item ige + -t / -tt képző
    \begin{itemize}
        \item a \underline{megtanult} ...
        \item az \underline{elképzelt} ...
    \end{itemize}
    \item Sokszor múlt idejű igével esik egybe az alakja
\end{itemize}
\subsection{Beálló melléknévi igenév}
\begin{itemize}
    \item ige + -andó / -endő képző
    \begin{itemize}
        \item pl. megvalósítandó, elérendő
    \end{itemize}
\end{itemize}

\section{Határozói igenév}
\begin{itemize}
    \item ige + -va / -ve / -ván / -vén képző
    \begin{itemize}
        \item futva, énekelve, futván, énekelvén
    \end{itemize}
\end{itemize}

\end{document}